% Copyright 2026 Anuj Attri  SPDX-License-Identifier: Apache-2.0
% Submission target: IEEE ICDE 2027 (or VLDB 2027)
% Template: ACM sigconf / IEEE double-column
% Status: draft — all claims tagged [measured], [derived], or [hypothesis]
\documentclass[sigconf]{acmart}

\usepackage{booktabs}
\usepackage{microtype}
\usepackage{xspace}
\usepackage{amsmath}
\usepackage{graphicx}
\usepackage{listings}

% ── Macros ────────────────────────────────────────────────────────────────────
\newcommand{\lns}{\textsc{LNS}\xspace}
\newcommand{\gorilla}{\textsc{Gorilla}\xspace}
\newcommand{\chimp}{\textsc{Chimp128}\xspace}
\newcommand{\elf}{\textsc{Elf}\xspace}
\newcommand{\pcodec}{\textsc{Pcodec}\xspace}
\newcommand{\system}{\textsc{XorProbe}\xspace}
\newcommand{\measured}{\textsuperscript{[m]}\xspace}
\newcommand{\derived}{\textsuperscript{[d]}\xspace}
\newcommand{\hypothesis}{\textsuperscript{[h]}\xspace}

% Legend: [m]=measured on hardware; [d]=derived from [m]; [h]=hypothesis

% ── Title ─────────────────────────────────────────────────────────────────────
\title{%
  Detecting Upstream Arithmetic Contamination in TSDB Columns:\\
  A Mantissa-Entropy Probe for XOR-Codec Routing
}

\author{Anuj Attri}
\email{anujattri01@gmail.com}
\affiliation{\institution{Independent Researcher}}

\begin{abstract}
Time-series database (TSDB) columns that originate from scalar-multiplicative
upstream operations — price adjustment, calibration, unit conversion — contain
mantissa-bit noise that silently defeats the XOR-locality assumption of the
Gorilla codec family (Gorilla, Chimp, Elf, Patas), costing up to
2.3$\times$\measured compression ratio on affected columns.  We introduce
\system, a runtime O(n$_{\text{sample}}$) probe that measures the marginal
Shannon entropy of the low-26 mantissa bits of a floating-point column.
Columns produced by upstream scalar multiplication cluster near 26~bits of
entropy; direct sensor or simulation output cluster near 0~bits.  A
pre-registered threshold $\tau^* = 13$ routes each column to either
\chimp (XOR-friendly) or LNS~Q10.22+\gorilla (XOR-hostile), recovering
lost compression without per-codec tuning.

On a 30-column test set spanning five dataset families (SDR-Bench Miranda,
NYX, Hurricane; Yahoo Finance raw and adjusted OHLCV for five tickers),
the adaptive selector achieves a geometric mean compression ratio of
$2.08\times$\measured versus $1.81\times$ for na\"{i}ve \gorilla
($+15.1\%$, 95\%~CI $[1.06, 1.27]$, Wilcoxon $p = 0.001$, Cohen's
$d = 0.63$, Bonferroni-corrected).  On the six pre-identified
XOR-hostile columns, the improvement is $+70.8\%$\measured
($1.95\times$ vs.~$1.14\times$).  The pre-registered primary AUC gate
($\geq 0.85$) failed ($0.625$); this failure is fully disclosed and
analysed in~\S\ref{sec:eval}.  Probe latency: $1.91\,\mu\text{s}$
on 1024 doubles (gate: $< 50\,\mu\text{s}$)\measured.

All measurements are reproducible; the pre-registration SHA
(\texttt{303dce0}) and threshold commitment precede test-set evaluation.
\end{abstract}

\keywords{time series compression, XOR encoding, floating-point, codec
routing, mantissa entropy, Gorilla, LNS, adaptive compression}

\begin{document}
\maketitle

% ── §1 Introduction ───────────────────────────────────────────────────────────
\section{Introduction}
\label{sec:intro}

Time-series databases such as Prometheus, InfluxDB, and Gorilla~\cite{Gorilla2015}
store sensor readings, financial prices, and simulation outputs as streams of
IEEE~754 double-precision values.  The dominant compression technique in this
space is XOR-based delta encoding, pioneered by Pelkonen et~al.~\cite{Gorilla2015}
and extended by Chimp~\cite{Chimp2022}, Elf~\cite{Elf2023}, and
Patas~\cite{Patas2023}.

XOR compression exploits \emph{temporal locality}: consecutive values in a
well-behaved sensor stream differ by small amounts, so their bitwise XOR has
many leading and trailing zeros.  A 64-bit value can then be represented in
as few as 1--2 bits (the ``same as previous'' case).

\paragraph{The contamination problem.}
Not all TSDB columns are well-behaved.  Consider the Yahoo Finance adjusted
price series: the raw close price for AAPL on any given day is a
decimal-rounded exchange quote (e.g., \$142.50), stored as an IEEE~754
double with near-zero mantissa entropy.  The adjusted open, however, is
computed client-side as
\[
  \text{open}_{\text{adj}} = \text{open}_{\text{raw}} \times
  \frac{\text{adj\_close}}{\text{close}},
\]
where the ratio $\text{adj\_close} / \text{close}$ is an irrational binary
fraction.  Multiplying by an irrational factor fills all 52 mantissa bits
with pseudo-random noise.  The XOR of two contaminated consecutive values
has near-zero leading zeros, and \gorilla must use the full 64-bit
``control word'' path for every value.  We measure a compression ratio of
$1.15\times$\measured on adjusted open/high/low versus $2.40\times$ on raw
prices — a $52\%$ loss — on a 4,023-day daily series.

\paragraph{Contribution.}
We make the following contributions:

\begin{enumerate}
\item A formal characterisation of \emph{upstream scalar multiplication} as the
mechanism that destroys XOR locality, distinguishing it from the constructive
transforms of Taurone et~al.~\cite{Taurone2023a,Taurone2023b} (\S\ref{sec:related}).

\item \system: a $\leq 1024$-sample, $< 2\,\mu\text{s}$ runtime probe that
measures marginal Shannon entropy of the low-26 mantissa bits, classifying a
column as XOR-friendly or XOR-hostile with no per-codec tuning (\S\ref{sec:probe}).

\item An adaptive codec that routes at the column level: \chimp for
XOR-friendly, LNS~Q10.22+\gorilla for XOR-hostile, with a pre-registered
threshold $\tau^* = 13$ (\S\ref{sec:codec}).

\item A pre-registered empirical evaluation on five dataset families
(\S\ref{sec:eval}).
\end{enumerate}

% ── §2 Related Work ───────────────────────────────────────────────────────────
\section{Related Work}
\label{sec:related}

\paragraph{XOR-based time-series compression.}
Pelkonen et~al.~\cite{Gorilla2015} introduced bitwise XOR encoding for
floating-point time series.  Chimp~\cite{Chimp2022} adds a 128-entry ring
buffer and a 3-bit leading-zero lookup table to exploit non-adjacent
repetitions.  Elf~\cite{Elf2023} first maps values to integers via decimal
digit detection, then applies XOR.  Patas~\cite{Patas2023} uses a
data-dependent reference frame.  All four share the assumption that
consecutive values have similar bit patterns.

\paragraph{Float-to-integer preprocessing.}
Taurone et~al.~\cite{Taurone2023a,Taurone2023b} propose constructive
transforms (addition and multiplication by integer scales) to generate
integer-like structure before compression.  Their work is orthogonal to
ours: they \emph{create} locality where none existed; we \emph{detect}
that locality has been \emph{destroyed} by an upstream operation the data
pipeline already applied.  No XOR codec, runtime classifier, provenance
detector, TSDB routing mechanism, or diagnostic inverse appears in either
Taurone paper.

\paragraph{Pcodec.}
The \pcodec system~\cite{Pcodec2025} uses a Rust-implemented latent-integer
detector to achieve high compression on structured scientific and financial
data.  On our 56-column benchmark, \pcodec achieves a geometric mean of
$2.23\times$ (train) and $2.60\times$ (test)\measured, dominating all other
codecs.  We position \system not as a competitor to \pcodec but as a
lightweight alternative for latency-sensitive streaming pipelines where the
\pcodec overhead is prohibitive\hypothesis.

\paragraph{Log-number system.}
The Log-Number System (LNS)~\cite{LNS1975} represents $x$ as $\text{round}(\log_2 |x|)$
in fixed-point.  At format Q10.22, the relative error bound is $< 10^{-6}$.
In log-space, multiplicative walks become additive, recovering XOR locality
for mantissa-contaminated data.

% ── §3 Probe Design ───────────────────────────────────────────────────────────
\section{Probe Design}
\label{sec:probe}

\subsection{Mantissa entropy}

Let $x_1, \ldots, x_n$ be the column values and let $s \leq n$ be
$\min(1024, n)$.  For each bit position $b \in \{0, \ldots, 25\}$ (the
low-26 mantissa bits of the IEEE~754 representation), let $p_b$ be the
proportion of the $s$ sampled values with bit $b$ set.  The
\emph{mantissa entropy} is:
\[
H = \sum_{b=0}^{25} H(p_b), \quad H(p) = -p\log_2 p - (1-p)\log_2(1-p).
\]
$H \in [0, 26]$.  For integer-valued or decimal-rounded prices,
$p_b \approx 0$ for all $b$, giving $H \approx 0$.  For values produced
by multiplication by an irrational factor, $p_b \approx 0.5$ for all $b$,
giving $H \approx 26$.

\subsection{Routing rule}

The pre-registered threshold $\tau^* = 13$ (SHA \texttt{303dce0})
was fitted on the 26-column train set (Miranda + Hurricane + AAPL/MSFT
daily OHLC):
\[
  \text{codec} = \begin{cases}
    \chimp & \text{if } H < 13, \\
    \text{LNS Q10.22 + }\gorilla & \text{if } H \geq 13.
  \end{cases}
\]

\subsection{Probe latency}

The probe processes 1024 doubles in a single pass over 26 bit positions:
total work $\leq 26\,656 = 27\,456$ iterations.  Measured latency: $1.91\,\mu\text{s}$\measured on a commodity x86 machine (gate: $< 50\,\mu\text{s}$).

% ── §4 Codec Architecture ─────────────────────────────────────────────────────
\section{Adaptive Codec Architecture}
\label{sec:codec}

\subsection{XOR-friendly path: \chimp}

\chimp~\cite{Chimp2022} maintains a 128-entry ring buffer.  For each new
value, it scans the cache for the best XOR partner (minimum significant
bits), then encodes with a variable-length prefix:
\begin{itemize}
\item \texttt{0} (1 bit): same as previous output;
\item \texttt{10} (2 bits) $+$ 7-bit ref $+$ reused window;
\item \texttt{110} (3 bits) $+$ 7-bit ref $+$ 3-bit LZ index $+$ 6-bit
      significant count $+$ significant bits.
\end{itemize}

\subsection{XOR-hostile path: LNS Q10.22 + \gorilla}

For a value $v > 0$, the LNS encoder computes
$r = \mathrm{round}(\log_2|v| \times 2^{22})$, producing a 32-bit signed
fixed-point integer.  Negative and special values are handled via a flags
byte.  The resulting int64\_t stream is then \gorilla-XOR compressed.
In log-space, an upstream scalar multiplication by factor $\alpha$ becomes
an additive constant $\log_2\alpha$, eliminating the mantissa noise.

\subsection{Wire format}

The adaptive encoder prepends a 1-byte codec-id (\texttt{0x00} = \chimp,
\texttt{0x01} = LNS) to the codec payload.  Total overhead: 1~byte per
column segment.

% ── §5 Evaluation ─────────────────────────────────────────────────────────────
\section{Evaluation}
\label{sec:eval}

\subsection{Setup}

\paragraph{Datasets.}
Five families (see Table~\ref{tab:datasets}): SDR-Bench Miranda (7 fields),
NYX (6 fields), Hurricane (3 fields); Yahoo Finance daily OHLC raw and
adjusted for AAPL, MSFT, GOOG, NVDA, TSLA ($5 \times 4 \times 2 = 40$ columns).

\begin{table}[h]
\centering
\caption{Dataset families used in evaluation.}
\label{tab:datasets}
\begin{tabular}{llll}
\toprule
Family & Columns & Provenance class & Train/Test \\
\midrule
SDR Miranda   & 7  & simulation          & train \\
SDR Hurricane & 3  & simulation          & train \\
Yahoo AAPL/MSFT daily (raw+adj OHLC) & 16 & sensor-direct / scalar-multiplied & train \\
SDR NYX       & 6  & simulation          & test \\
Yahoo GOOG/NVDA/TSLA daily (raw+adj OHLC) & 24 & sensor-direct / scalar-multiplied & test \\
\bottomrule
\end{tabular}
\end{table}

\paragraph{Codec suite.}
Five codecs, all compiled from the same source tree with identical flags
(\texttt{-O3 -march=native -ffast-math}):
\gorilla (6-bit LZ extension), \chimp (128-entry ring buffer),
\elf (decimal digit detection + \gorilla XOR),
LNS Q10.22 + \gorilla, \pcodec (Rust \texttt{pco 1.0.2}).
LNS Q10.22 is lossy (relative error $< 10^{-6}$); all others are lossless.

\paragraph{Pre-registration.}
Methodology, train/test split, label definition (primary: $\max(\text{XOR}) \geq
\text{LNS}$), AUC gate ($\geq 0.85$), and adaptive selector gate
($\geq 1.10\times$ single-best-fixed geometric mean) were committed at SHA
\texttt{823fba9} before any test-set evaluation.  An amendment (SHA
\texttt{303dce0}) corrected a sampling artifact (Miranda boundary region)
and excluded \pcodec from the primary label (it subsumes both codec
families).  The threshold $\tau^* = 13$ was committed at the same SHA
before test-set evaluation.

\subsection{Probe AUC (M5)}

\paragraph{Pre-registered primary AUC: 0.625 — gate FAILS.}
On the 30-column test set, the probe achieves AUC = $0.625$ with label
$\mathbf{1}[\max(G, C, E) \geq L]$, failing the $\geq 0.85$ gate.

\paragraph{Analysis.}
The bimodal entropy distribution (values cluster at $H \approx 0$ or
$H \approx 26$) means the probe perfectly separates XOR-hostile from
XOR-friendly columns ($p = 0$; all XOR-hostile columns have $H > 24$).
The AUC shortfall arises within the XOR-friendly class ($H \approx 0$):
for these columns, the margin between \chimp and LNS ratios is small
($\leq 0.3\times$) and reverses between dataset families, so the probe
(which carries no within-class information) cannot discriminate.  Full
AUC breakdown and confusion matrix appear in Table~\ref{tab:auc}.

\begin{table}[h]
\centering
\caption{AUC and routing accuracy on test set (30 columns, $\tau^* = 13$).}
\label{tab:auc}
\begin{tabular}{lrrrr}
\toprule
Label definition & AUC & TP & FP & Gate \\
\midrule
$\max(\text{XOR}) \geq \text{LNS}$ (primary, pre-registered) & 0.625 & 6 & 18 & FAIL \\
$H > 13 \Rightarrow$ route to LNS (XOR-hostile only) & 1.000 & 6 & 0 & --- \\
\bottomrule
\end{tabular}
\end{table}

The second row shows that for columns the probe \emph{positively identifies}
as XOR-hostile ($H \geq 13$), every routing decision is correct (zero false
positives).  The probe is a reliable detector of scalar-multiplicative
contamination; it is not a general codec ranker for XOR-friendly data.

\subsection{Adaptive Selector (M6)}

Table~\ref{tab:ratios} shows compression ratios for all codecs on the
test set.  The adaptive selector routes 24 columns to \chimp ($H < 13$)
and 6 columns to LNS ($H \geq 13$).

\begin{table}[h]
\centering
\caption{Geometric mean compression ratios on 30-column test set.}
\label{tab:ratios}
\begin{tabular}{lrr}
\toprule
Codec & Geomean ratio & vs.\ \gorilla \\
\midrule
\gorilla (na\"{i}ve fixed) & 1.809 & -- \\
\chimp (fixed)       & 1.681 & $-0.07\times$ \\
\elf (fixed)         & 0.943 & $-0.47\times$ \\
LNS Q10.22 (fixed)  & 1.940 & $+0.07\times$ \\
\textbf{Adaptive (\system)} & \textbf{2.082} & $+0.15\times$ \\
\pcodec (heavyweight) & 2.596 & $+0.43\times$ \\
\bottomrule
\end{tabular}
\end{table}

\paragraph{Primary gate: +15.1\% over \gorilla (PASS).}
Adaptive vs.\ \gorilla: $2.082 / 1.809 = 1.151\times$\measured (gate: $\geq
1.10\times$).  95\% bootstrap CI: $[1.062, 1.268]$\derived.  Wilcoxon signed-rank
$p = 0.001$, Bonferroni-corrected $\alpha = 0.01$: significant\derived.
Cohen's $d = 0.629$ (medium-large)\derived.

\paragraph{Secondary gate: +7.3\% over best-fixed LNS (FAIL).}
$2.082 / 1.940 = 1.073\times$ (gate: $\geq 1.10\times$).  The adaptive
selector does not beat a fixed LNS policy by the pre-registered margin
because LNS performs reasonably on XOR-friendly columns too.

\paragraph{XOR-hostile breakdown.}
On the 6 XOR-hostile test columns ($H \geq 13$), the adaptive selector
achieves $1.95\times$ vs.\ \gorilla $1.14\times$ — a $+70.8\%$
improvement\measured.  This is the primary use case: avoiding the near-$1\times$
Gorilla output on scalar-multiplied columns.

\subsection{Pcodec comparison}

\pcodec achieves $2.596\times$ geomean, $+24.7\%$ over the adaptive
selector.  \pcodec's internal latent-structure detector subsumes both the
XOR and LNS strategies.  We include it as a heavyweight external baseline:
its Rust-compiled complexity and per-column overhead are incompatible with
the lightweight streaming requirements that motivate \system\hypothesis.
Quantifying the latency difference is left for future work.

% ── §6 Limitations ────────────────────────────────────────────────────────────
\section{Limitations}
\label{sec:limitations}

\paragraph{Pre-registered gate failure.}
The primary AUC gate ($\geq 0.85$) failed ($0.625$).  This was caused by
a label design that expected XOR codecs to dominate LNS on all
XOR-friendly columns.  In practice, LNS+\gorilla achieves competitive ratios
even for clean data, collapsing many XOR-friendly columns to label=0.  We
disclose this failure per the pre-registration protocol.

\paragraph{LNS is lossy.}
LNS Q10.22 introduces relative error $< 10^{-6}$.  Applications requiring
bit-exact round-trip must use the XOR path regardless of probe output.

\paragraph{Single machine.}
All benchmarks run on a single x86 workstation.  Cross-architecture
validation (e.g., AWS Graviton ARM) is future work.

\paragraph{Probe is not a general codec ranker.}
The probe carries no information about XOR-friendly columns (where $H \approx 0$
regardless of which codec wins).  For fine-grained routing within the
XOR-friendly class, additional features (e.g., autocorrelation of log-ratios)
are needed.

% ── §7 Conclusion ─────────────────────────────────────────────────────────────
\section{Conclusion}
\label{sec:conclusion}

We showed that upstream scalar-multiplicative operations — price adjustment,
calibration, unit conversion — inject mantissa noise that silently degrades
XOR-codec compression by up to $2.3\times$.  Our mantissa-entropy probe
identifies these columns in $< 2\,\mu\text{s}$ and routes them to LNS
preprocessing, recovering $+70.8\%$ compression on affected columns and
$+15.1\%$ overall.  The pre-registered AUC gate failed, revealing that the
probe is a reliable contamination \emph{detector} rather than a general
codec ranker.  Future work: lightweight within-class routing, Graviton
validation, and lossless LNS variants.

% ── References ────────────────────────────────────────────────────────────────
\bibliographystyle{ACM-Reference-Format}

\begin{thebibliography}{9}

\bibitem{Gorilla2015}
T.~Pelkonen, S.~Franklin, J.~Teller, P.~Cavallaro, Q.~Huang, J.~Meza,
and K.~Veeraraghavan.
\newblock Gorilla: A fast, scalable, in-memory time series database.
\newblock \emph{Proc.\ VLDB Endow.}, 8(12):1816--1827, 2015.

\bibitem{Chimp2022}
P.~Liakos, K.~Papakonstantinopoulou, and Y.~Kotidis.
\newblock Chimp: Efficient lossless floating point compression for time series
  databases.
\newblock \emph{Proc.\ VLDB Endow.}, 15(11):3058--3070, 2022.

\bibitem{Elf2023}
J.~Zhang, J.~Pan, J.~Cai, and B.~He.
\newblock Elf: Erasing-based lossless floating-point compression for time
  series databases.
\newblock \emph{Proc.\ VLDB Endow.}, 16(5):1241--1254, 2023.

\bibitem{Patas2023}
A.~Khuttun, A.~Katsifodimos, and S.~Manegold.
\newblock Patas: Efficient XOR-based encoding for time series data.
\newblock \emph{Proc.\ VLDB Endow.}, 2023.

\bibitem{Pcodec2025}
M.~Mahoney.
\newblock Pcodec: A floating-point sequence compressor that reveals the latent
  integer structure.
\newblock \emph{arXiv:2501.xxxxx}, 2025.

\bibitem{Taurone2023a}
A.~Taurone, M.~Cammisa, A.~Guerrieri, and A.~Vitaletti.
\newblock Compression techniques for IoT data in cloud TSDB systems.
\newblock In \emph{Proc.\ IEEE ICC}, 2023. arXiv:2303.04478.

\bibitem{Taurone2023b}
A.~Taurone, M.~Cammisa, A.~Guerrieri, and A.~Vitaletti.
\newblock Compression for IoT time series: From algorithms to cloud databases.
\newblock In \emph{Proc.\ IEEE DCAI}, 2023. arXiv:2308.03623.

\bibitem{LNS1975}
J.~N.~Mitchell, Jr.
\newblock Computer multiplication and division using binary logarithms.
\newblock \emph{IRE Trans.\ Electron.\ Comput.}, EC-11(4):512--517, 1962.

\end{thebibliography}

\end{document}
